Jusqu'à présent nous avons tracé nos dessins sur du papier quadrillé à petits carreaux.
Le \robot\ a la faculté supplémentaire de changer la taille des carreaux. 
On peut donc agrandir les dessins en augmentant la taille des carreaux avant de commencer le dessin. Il est également possible 
de changer la taille des carreaux pendant que le \robot\ dessine. Ceci permet de
tracer des segments de longueurs différentes et des arcs de cercles dont le rayon peut varier. 
En choisissant des carreaux rectangulaires on peut tracer des ellipses et choisir
de manière quelconque les angles de déplacement du \robot.

\section{Grille carrée}
La grille standard représentée à la figure \ref{fig:grille} (page \pageref{fig:grille})
a une taille de 12 pixels. On peut changer la 
\underline{l}ongueur du côté d'un carreau de la \underline{g}rille avec la commande {\idt lg} qui
prend en paramètre la nouvelle longueur de la grille exprimée en pixels. Par exemple
{\idt lg(6)} définit une taille de grille deux fois plus petite que la grille habituelle tandis que 
{\idt lg(24)} double la taille de la grille.

La longueur de la grille peut aussi
comporter une partie décimale, comme dans {\idt lg(7.5)} (7 pixels et demi). Un tel
nombre, que l'on appelle un {\em nombre flottant}\index{Nombre flottant}, s'écrit avec un point
(et non pas avec une virgule comme en français)~:
{\idt 6.0} (ou {\idt 6.}) est le nombre flottant six
tandis que {\idt 6} est le nombre entier six. Les commandes {\idt lg}, {\idt lgX} et {\idt lgY}
acceptent indifféremment un nombre entier ou un nombre flottant~: {\idt lg(6)} et {\idt lg(6.0)}
ont le même effet\footnote{\ Ce n'est pas le cas de tous les langages~: en \ocaml\ par exemple,
il faut obligatoirement écrire {\idt lg\ 6.} avec un nombre flottant et convertir les entiers
en flottants.}.
Quand la longueur de la grille est donnée par une expression entière comme 
$(6\ \ast\ ${\idt i}$)$, on écrit donc simplement {\idt lg(6\ $\ast$\ i)}. Nous étudierons
les nombres flottants en détail au chapitre qui leur est consacré.

\begin{exemple}[Carrés superposés]{\rm
 Le programme ci-dessous permet de des\-si\-ner des carrés superposés, le \robot\ partant
 et arrivant au coin inférieur gauche, orienté vers le nord, la superposition étant
 obtenue en choisissant des tailles de grille successives de 6, 12, 18, \ldots, 84~:
 
% \includePostScript{CarresSuperposes2}{0.95}{0.95}
% incorpore dans~:
 \input python/CarresSuperposes2.tex
\noindent $\Box$ }\end{exemple}


\section{Grille rectangulaire}

Le \robot\ peut se promener sur une grille rectangulaire dont la longueur du côté
horizontal est différente de celle du côté
vertical. Ceci lui per\-met de des\-sin\-er des rec\-tan\-gles et des el\-lip\-ses com\-me ceux-ci~:
%de la figure \ref{fig:RecEl} (page \pageref{fig:RecEl}).
\begin{figure}[htbp]
 \includevspacePostScript[2.2cm]{RectangleEtEllipse40x20}{0.95}{0.95}
 \caption{Rectangle et ellipse sur un quadrillage $40 \times 20$}
% \label{fig:RecEl}
\end{figure}

 La grille initiale a une taille $12 \times 12$ comptée en pixels. La commande {\idt lgX(40)} permet de donner
 une longueur de 40 pixels au côté horizontal de la grille sans changer la longueur du côté vertical. après
 cette commande la grille a donc une taille $40 \times 12$ pixels. Dans cette notation $40 \times 12$ on donne
 d'abord la longueur du côté horizontal (nommé $X$) puis la longueur du côté vertical (nommé $Y$)
 de la grille. La commande {\idt lgY(20)} permet de donner
 une longueur de 20 pixels au côté vertical de la grille sans changer la longueur du côté horizontal. après
 cette commande la grille a donc une taille $40 \times 20$ pixels. On peut donc changer indépendamment la
  taille horizontale (avec {\idt lgX}) ou la taille verticale (avec {\idt lgY}) de la grille. 
  On peut revenir à une grille carrée en utilisant la commande {\idt lg}.
  Par exemple la commande {\idt lg(12)}
  permet de revenir à une grille carrée standard $12 \times 12$. Le rectangle et l'ellipse
 ci-dessus 
  %de la figure \ref{fig:RecEl} (page \pageref{fig:RecEl}) 
  sont construits par le programme suivant~:\hfill
 \input python/RectangleEtEllipse.tex
 
 La figure 
 ci-dessous
%\ref{fig:vdavvg3060} (page \pageref{fig:vdavvg3060}) 
illustre l'effet des commandes
 {\idt vd()},  {\idt av()} et {\idt vg()} avec une grille rectangulaire de dimensions $30 \times 60$.
\begin{figure}[htbp]
 \includevspacePostScript[4.3cm]{VD+AV+VG-30x60}{0.95}{0.95}
 \caption{Virage  à droite, avancer et virage à gauche avec une grille $30 \times 60$}
% \label{fig:vdavvg3060}
\end{figure}
\begin{figure}[htbp]
 \includevspacePostScript[2.9cm]{Parallel+Ellipses}{0.95}{0.95}
 \caption{Parallélogrammes et ellipses}
 \label{fig:parel}
\end{figure}
\begin{exemple}{\rm 
\renewcommand{\carrenoir}{\hspace{-4mm}\rule{1.5mm}{1.5mm}\hspace{2.5mm}}
Les divers tracés de segments et d'arcs d'el\-lip\-ses par les commandes {\idt av()}, {\idt vg()} 
et {\idt vd()} sont illustrés par les deux programmes ci-dessous qui dessinent des parallélogrammes et
des ellipses par pivotement d'un quart de tour autour du centre de la grille, comme on le voit à la 
fi\-gu\-re \ref{fig:parel} (page \pageref{fig:parel})~:\hfill
\input python/Parallelogrammes.tex
\input python/Ellipses.tex
\renewcommand{\carrenoir}{}
\noindent $\Box$}\end{exemple}

\noindent \rule{137mm}{0.1mm} %---------------------------------------------
\begin{exercices}{Utiliser les possibilités offertes par les changements de taille de la grille carrée
pour dessiner~:

\tiret les spirales et cercles concentriques de la figure \ref{fig:spirales} (page \pageref{fig:spirales})$\:$;

\tiret la frise de la figure \ref{fig:dfrisec} (page \pageref{fig:dfrisec})$\:$;

\tiret l'hélice et la balle de la figure \ref{fig:heba} (page \pageref{fig:heba})$\:$;

\tiret l'escargot de la figure \ref{fig:escar} (page \pageref{fig:escar})$\:$;

\tiret le stégosaure et son bébé de la figure \ref{fig:steg2} (page \pageref{fig:steg2})$\:$;

\tiret la famille de chiens de la figure \ref{fig:chiens} (page \pageref{fig:chiens})$\:$;

\tiret le pavage de la figure \ref{fig:PVdV} (page \pageref{fig:PVdV}) inspiré par un motif de
van de Vecht (\cite{escher1}, page 26) représenté à la figure \ref{fig:mVdV} (page \pageref{fig:mVdV}).
\noindent $\Box$}\end{exercices}
\begin{exercices}{Utiliser les possibilités offertes par les  grilles rectangulaires
pour dessiner~:

\tiret la frise de losanges se chevauchant de la figure \ref{fig:frlos} (page \pageref{fig:frlos})$\:$;

\tiret le vélo de la figure \ref{fig:velo} (page \pageref{fig:velo})$\:$;

\tiret la rose des vents
à huit directions \ref{fig:rose-des-vents-8d} (page \pageref{fig:rose-des-vents-8d})$\:$;

\tiret les pavages des figures \ref{fig:PVdV1l} (page \pageref{fig:PVdV1l}),
\ref{fig:PVdV2l} (page \pageref{fig:PVdV2l}) et \ref{fig:PVdV4l} (page \pageref{fig:PVdV4l}) obtenues à
partir du pavage de la figure \ref{fig:PVdV} (page \pageref{fig:PVdV}) en diminuant pro\-gres\-si\-ve\-ment la taille de la grille sur le côté droit
du dessin, sur le côté gauche et dans le dernier cas sur les quatre côtés du rectangle à la manière du graveur
hollandais M. C. Escher (\cite{escher1}, \cite{escher2}).

\noindent $\Box$}\end{exercices}
\begin{figure}[htbp]
 \includevspacePostScript[2.8cm]{Spirales}{0.95}{0.95}
 \caption{Spirales et cercles concentriques}
 \label{fig:spirales}
\end{figure}
\begin{figure}[htbp]
 \includevspacePostScript[0.6cm]{DoubleFriseCrantee}{0.95}{0.95}
 \caption{Double frise crantée}
 \label{fig:dfrisec}
\end{figure}
\begin{figure}[htbp]
 \includevspacePostScript[4.3cm]{HeliceBalle}{0.95}{0.95}
 \caption{Hélice et balle}
 \label{fig:heba}
\end{figure}
\begin{figure}[htbp]
 \includevspacePostScript[4.8cm]{Escargot}{0.95}{0.95}
 \caption{Escargot}
 \label{fig:escar}
\end{figure}
\begin{figure}[htbp]
 \includevspacePostScript[5.8cm]{StegosaureEtBebe}{0.95}{0.95}
 \caption{Stégosaure et son bébé}
 \label{fig:steg2}
\end{figure}
\begin{figure}[htbp]
 \includevspacePostScript[4.8cm]{ChiensEtChiots}{0.95}{0.95}
 \caption{Famille de chiens en balade}
 \label{fig:chiens}
\end{figure}
\begin{figure}[htbp]
 \includevspacePostScript[2.5cm]{MotifVanDeVecht}{0.95}{0.95}
 \caption{Motif ornemental de N. J. van de Vecht}
 \label{fig:mVdV}
\end{figure}
\begin{figure}[htbp]
 \includevspacePostScript[2.2cm]{FriseLosangesImbriques}{0.95}{0.95}
 \caption{Frise de losanges se chevauchant}
 \label{fig:frlos}
\end{figure}
\begin{figure}[htbp]
 \includevspacePostScript[4.1cm]{velo}{0.95}{0.95}
 \caption{Vélo}
 \label{fig:velo}
\end{figure}
\begin{figure}[htbp]
 \includevspacePostScript[14.6cm]{PavageVdVRegulier}{0.95}{0.95}
 \caption{Pavage régulier inspiré par un motif ornemental\label{fig:PVdV}}
 \centering{de N. J. van de Vecht}
\end{figure}
\begin{figure}[htbp]
 \includevspacePostScript[14.6cm]{PavageVdV1limite}{0.95}{0.95}
 \caption{Pavage rectangulaire à une limite}
 \label{fig:PVdV1l}
\end{figure}
\begin{figure}[htbp]
 \includevspacePostScript[14.6cm]{PavageVdV2limites}{0.95}{0.95}
 \caption{Pavage rectangulaire à deux limites}
 \label{fig:PVdV2l}
\end{figure}
\begin{figure}[htbp]
 \includevspacePostScript[17.0cm]{PavageVdV4limites}{0.95}{0.95}
 \caption{Pavage rectangulaire à quatre limites}
 \label{fig:PVdV4l}
\end{figure}

\clearpage
%%%%%%%%%%%%%%%%%%%%%%%%%%%%%%%%%%%%%%%%%%%%%%%%%%%
\begin{corrige}{
%\label{cor:Frise-1}
\renewcommand{\carrenoir}{\hspace{-4mm}\rule{1.5mm}{1.5mm}\hspace{2.5mm}}
Nous donnons quelques programmes, les autres programmes se trouvant dans le dossier
{\idt livre/\allowbreak 10\_grille\_cartesienne}~:
\input python/SpiraleGeometrique.tex
\input python/Escargot2.tex
\renewcommand{\carrenoir}{}
\noindent $\Box$}\end{corrige}
\begin{corrige}{
Nous donnons le programme le plus difficile, à savoir le pavage de la figure
\ref{fig:PVdV4l} (page \pageref{fig:PVdV4l}) obtenue à
partir du pavage de la figure \ref{fig:PVdV} (page \pageref{fig:PVdV}) en diminuant pro\-gres\-si\-ve\-ment la tail\-le de la gril\-le sur les 
qua\-tre côtés du rectangle. Les autres programmes se trouvent
dans le dossier {\idt livre/\allowbreak 10\_grille\_cartesienne}~:
\input python/PavageVdV4limites.tex
\noindent $\Box$}\end{corrige}

%%%%%%%%%%%%%%%%%%%%%%%%%%%%%%%%%%%%%%%%%%%%%%%%%%%%%%%%%%%%%%%%%%%%%%%%%%%
\noindent \framebox[\textwidth]{
\begin{minipage}{\myparboxwidth}
\setlength{\parindent}{2em}
\small \sf
\vspace{\abovedisplayskip}
\section*{Addenda}
Si toutes les grilles des points d'arrêt du \robot\ étaient alignées sur le centre de l'écran, il se pourrait
qu'après un changement de grille le \robot\ ne soit plus sur un point de la grille. C'est pourquoi la
 grille est toujours relative au \robot\ et non pas à l'écran. En cas de changement de taille de la grille tous les points de la grille (que l'on peut observer grâce à
la commande {\idt dg()}) changent de place, sauf celui sur lequel est placé le \robot. En conséquence aucune commande
du \robot\ (sauf une, à savoir {\idt ce()}) ne peut lui faire quitter l'un des points de la grille courante. Par exemple les commandes {\idt ag()}, {\idt ad()}, {\idt eh()} et 
{\idt eb()} déplacent le \robot\ horizontalement ou verticalement sur le point de la grille courante (respectivement à gauche, à droite, en haut et
en bas) le plus près du bord. 
Pour pouvoir placer le \robot\ n'importe où sur l'écran avec précision, il faut disposer d'un
 point de repère absolu. Ce point de repère absolu est le centre de l'écran. La commande 
 {\idt ce()} ramène le \robot\ toujours au même point exactement situé au centre de l'écran. 
 Cette commande est donc la seule
 commande du \robot\ qui ne s'effectue pas sur la même grille.  Si le centre
 de l'écran n'est pas situé sur la grille courante, la commande {\idt ce()} recentre la grille exactement 
 au centre de l'écran de façon à éviter de faire sortir le \robot\ de son quadrillage de déplacement.

Pour éviter d'éventuels problèmes de décalage, il suffit de programmer le dessin comme s'il s'agissait
de papier millimétré en  choisissant une taille de grille minimale (1 millimètre sur le papier  millimétré)
de façon  à ce que tous les changements de taille de grille soient effectués quand le \robot\ est au centre et correspondent
à des multiples (5, 10, \ldots sur le papier  millimétré) de la taille de base. On peut également procéder
empiriquement par expérimentation, ce qui est quelquefois plus distrayant.

\vspace{\abovedisplayskip}
\end{minipage}
}
%%%%%%%%%%%%%%%%%%%%%%%%%%%%%%%%%%%%%%%%%%%%%%%%%%%%%%%%%%%%%%%%%%%%%%%%%%%%
